\scanpage{30}\mspage{24 bis}
% printed form: blank music staves.
% @ top, heading | text | Position relative d'un triangle
\uline{Position relative d'un triangle $t$ et d'une base $b$}.

% @ top, root pairing | text | On a
On a $\sum_{i\in t}i=\lambda$ orthogonal aux racines, donc
% @ top, sum | display | Σ i.r_b = 0
\[
\sum_{i\in t}i.r_b=0
\]
% @ upper third, alternatives | text | et comme i r_b
et comme $i\,r_b\in\{0,1,-1\}$ suivant la position relative de $r_b$ et $i$ i.e. \inserted{de} $b$ et $i$, on trouve 2 possibilités (à conjugaison près)

% @ upper third, a | text | i r_b = 0
a) $i\,r_b=0$ $\forall i\in t$, i.e. \deleted{\unsure{de|$E$}} $r_b$ (ou $b$) orthogonal à $t$ : déjà étudié. L'ens de ces $b$ est au nb de 24, stable par passage à l'opposée.
\margin{Ainsi, les bases orthogonales à $t$ sont celles qui sont disjointes de $t$.}% slanted

% @ upper middle, b | text | r_b non orthogonal
b) \deleted{\unsure{\#|$\exists$}} $r_b$ non orthogonal à $t$. Alors on peut écrire $t=\{u,v,w\}$ avec
% @ upper middle, brace at right | display | r_b(u), r_b(v), r_b(w)
\[
\left\{\begin{aligned}
r_b(u)&=1\quad\text{i.e. }u\in b\ \text{i.e. }E(u)\cap b=\varnothing\\
r_b(v)&=-1\quad\text{i.e. }\unsure{u|v}\in b'\ \text{i.e. }\unsure{\widehat E(v)|E(v)}\cap b\ \text{de card. }5\\
r_b(w)&=0\quad\text{i.e. }w\notin b\cup b'\ \text{i.e. }E(w)\cap b\ \text{card }2
\end{aligned}\right.
\]
% @ middle, count and question | text | Pour t fixé
Pour $t$ fixé, le nb de ces $b$ est $72-24=48$, se groupant en 6 groupes de 8 bases $b$, pour lesquels $u,v,w$ sont fixés. Comment paramétrer ces 8 bases ? On note que \inserted{l'intersection de} $E(w)$ avec $b$ et $b'$ est de cart. 2, il contient exactement un pt $u'\in b'$ et un $v'\in b$, et ces deux sont liés. Donc les \unsure{8} choix s'obtiennent en choisissant \deleted{1°) le triangle \unsure{$s$|$S$|6} $(w,u',v')$ parmi les quatre triangles $\ne(w,v,u)$ ayant $w$ comme sommet 2°) un des huit sommets de \inserted{$E(w)^*$} \unsure{5 distincts de $w$ comme étant $u'$|5 distincts de $u$ comme étant $u'$},} la base $b$ est l'unique base \inserted{de $C$} pour laquelle $u\in b$, $u'\in b'$,

% @ foot, reformulation | text | Le choix d'un couple
Le choix d'un couple $(t,b)$ en position non orthogonale équivaut donc à celle d'un triangle épinglé $\{u,v,w\}$ et d'un $u'\in E(w)^*$ (Nb : $27.10.1.8$)
% archive sheet number: 29
